\subsubsection{Linear Systems of ODEs and Asymptotic Stability}\label{subsec:linear-asymptotic-stability}

We now examine a concept called asymptotic stability in the context of linear system of ODEs. We consider the problem of finding a function $\bm{y}:\mathbb{R}\to\mathbb{R}^n$ that satisfies
\begin{equation}\label{eq:linear-nonhomogeneous-system}
\dfrac{\mathrm{d}\bm{y}}{\mathrm{d}t}(t)=\bm{A}(t)\bm{y}(t)+\bm{f}(t),
\end{equation}

\begin{equation}\label{eq:linear-system-initial-data}
\bm{y}(t_0)=\bm{y}_0,
\end{equation}
where $t_0\in\mathbb{R}$, $\bm{y}_0\in\mathbb{R}^n$, the vector valued function $\bm{f}:\mathbb{R}\to\mathbb{R}^n$ and the matrix valued function $\bm{A}:\mathbb{R}\to\mathbb{R}^{n\times n}$ are given.

Asymptotic stability describes the behavior of solutions of homogeneous systems as $t$ goes to infinity.

\begin{Definition}
\label{def:linear-system-asymptotic-stability}
The linear homogeneous system
\begin{equation}\label{eq:linear-homogeneous-system}
\bm{y}'=\bm{A}(t)\bm{y}
\end{equation}
is
\begin{enumerate}
\item \emph{asymptotically stable} if every solution of \eqref{eq:linear-homogeneous-system} satisfies
\begin{equation}\label{eq:asymptotic-stability-limit}
\lim_{t\to\infty}|\bm{y}(t)|=0,
\end{equation}

\item \emph{completely unstable} if every nonzero solution of \eqref{eq:linear-homogeneous-system} satisfies
\begin{equation}\label{eq:complete-instability-limit}
\lim_{t\to\infty}|\bm{y}(t)|=\infty.
\end{equation}
\end{enumerate}
\end{Definition}

The following fundamental result applies to constant coefficient systems.

\begin{Theorem}
\label{thm:constant-coefficient-stability}
Let $\bm{A}\in\mathbb{R}^{n\times n}$ be a constant matrix with eigenvalues
\[
\lambda_1,\lambda_2,\ldots,\lambda_n.
\]
Then the linear homogeneous system of ODEs
\begin{equation}\label{eq:constant-coefficient-homogeneous-system}
\bm{y}'=\bm{A}\bm{y}
\end{equation}
is
\begin{enumerate}
\item asymptotically stable if and only if all the eigenvalues of $\bm{A}$ have negative real parts; and
\item completely unstable if and only if all the eigenvalues of $\bm{A}$ have positive real parts.
\end{enumerate}
\end{Theorem}

The proof of this theorem is based on a diagonalization procedure for the matrix $\bm{A}$ and the following formula for all solutions of the initial-value problem associated with \eqref{eq:constant-coefficient-homogeneous-system}
\begin{equation}\label{eq:matrix-exponential-solution}
\bm{y}(t):=e^{\bm{A}(t-t_0)}\bm{y}_0.
\end{equation}
Here the matrix $e^{\bm{A}t}$ is defined by the uniformly convergent power series
\begin{equation}\label{eq:matrix-exponential-series}
e^{\bm{A}t}:=\sum_{n=0}^{\infty}\dfrac{\bm{A}^nt^n}{n!}.
\end{equation}


Formula \eqref{eq:matrix-exponential-solution} is the precursor of formulas in semigroup theory that we encounter in Chapter~\ref{chap:semigroup-methods}.

\subsubsection{Well-Posed Problems}\label{subsec:wellposed-problems}

We say that a problem is \emph{well-posed (in the sense of Hadamard)} if
\begin{enumerate}
\item there exists a solution,
\item the solution is unique,
\item the solution depends continuously on the data.
\end{enumerate}

If these conditions do not hold, a problem is said to be \emph{ill-posed}. Of course, the meaning of the term \emph{continuity with respect to the data} has to be made more precise by a choice of norms in the context of each problem considered.

In the course of this book we classify most of the problems we encounter as either well-posed or ill-posed, but the reader should avoid the assumption that well-posed problems are always "better" or more "physically realistic" than ill-posed problems. As we saw in the problem of buckling of a beam mentioned above, there are times when the conditions of a well-posed problem (uniqueness in this case) are physically unrealistic. The importance of ill-posedness in nature was stressed long ago by Maxwell~\cite{Max}:

\begin{quote}
For example, the rock loosed by frost and balanced on a singular point of the mountain-side, the little spark which kindles the great forest, the little word which sets the world afighting, the little scruple which prevents a man from doing his will, the little spore which blights all the potatoes, the little gemmule which makes us philosophers or idiots. Every existence above a certain rank has its singular points: the higher the rank, the more of them. At these points, influences whose physical magnitude is too small to be taken account of by a finite being may produce results of the greatest importance. All great results produced by human endeavour depend on taking advantage of these singular states when they occur.
\end{quote}

We draw attention to the fact that this statement was made a full century before people "discovered" all the marvelous things that can be done with cubic surfaces in $\mathbb{R}^3$.

\subsubsection{Representations}\label{subsec:representations}

There is one way of proving existence of a solution to a problem that is more satisfactory than all others: writing the solution explicitly. In addition to the aesthetic advantages provided by a representation for a solution there are many practical advantages. One can compute, graph, observe, estimate, manipulate and modify the solution by using the formula. We examine below some representations for solutions that are often useful in the study of PDEs.

\paragraph*{Variation of parameters}

Variation of parameters is a formula giving the solution of a nonhomogeneous linear system of ODEs \eqref{eq:linear-nonhomogeneous-system} in terms of solutions of the homogeneous problem \eqref{eq:linear-homogeneous-system}. Although this representation has at least some utility in terms of actually computing solutions, its primary use is analytical.

The key to the variations of constants formula is the construction of a \emph{fundamental solution matrix} $\Phi(t,\tau)\in\mathbb{R}^{n\times n}$ for the linear homogeneous system. This solution matrix satisfies
\begin{equation}\label{eq:fundamental-matrix-dynamics}
\dfrac{\mathrm{d}}{\mathrm{d}t}\Phi(t,\tau)=\bm{A}(t)\Phi(t,\tau),
\end{equation}

\begin{equation}\label{eq:fundamental-matrix-normalization}
\Phi(\tau,\tau)=\bm{I},
\end{equation}
where $\bm{I}$ is the $n\times n$ identity matrix. The proof of existence of the fundamental matrix is standard and is left as an exercise.

Note that the unique solution of the initial-value problem \eqref{eq:linear-homogeneous-system}, \eqref{eq:linear-system-initial-data} for the homogeneous system is given by
\begin{equation}\label{eq:fundamental-matrix-homogeneous-solution}
\bm{y}(t):=\Phi(t,t_0)\bm{y}_0.
\end{equation}


The use of Leibniz' formula reveals that the \emph{variation of parameters formula}
\begin{equation}\label{eq:variation-of-parameters}
\bm{y}(t):=\Phi(t,t_0)\bm{y}_0+\int_{t_0}^t\Phi(t,s)\bm{f}(s)\,\mathrm{d}s
\end{equation}
gives the solution of the initial-value problem \eqref{eq:linear-nonhomogeneous-system}, \eqref{eq:linear-system-initial-data} for the nonhomogeneous system.

\paragraph*{Cauchy's integral formula}

Cauchy's integral formula is the most important result in the theory of complex variables. It provides a representation for analytic functions in terms of its values at distant points. Note that this representation is rarely used to actually compute the values of an analytic function; rather it is used to deduce a variety of theoretical results.

\begin{Theorem}{Cauchy's integral formula}
\label{thm:cauchy-integral-formula}
Let $f$ be analytic in a simply connected domain $D\subset\mathbb{C}$ and let $C$ be a simple closed positively oriented curve in $D$. Then for any point $z_0$ in the interior of $C$
\begin{equation}\label{eq:cauchy-integral-formula}
f(z_0)=\dfrac{1}{2\pi i}\oint_C\dfrac{f(z)}{z-z_0}\,\mathrm{d}z.
\end{equation}

\end{Theorem}

\subsubsection{Estimation}\label{subsec:estimation}

When we speak of an \emph{estimate} for a solution we refer to a relation that gives an indication of the solution's size or character. Most often these are inequalities involving norms of the solution. We distinguish between the following two types of estimate. An \emph{a posteriori} estimate depends on knowledge of the existence of a solution. This knowledge is usually obtained through some sort of construction or explicit representation. An \emph{a priori} estimate is one that is conditional on the existence of the solution; i.e., a result of the form, "If a solution of the problem exists, then it satisfies ..." We present here an example of each type of estimate.

\paragraph*{Gronwall's inequality and energy estimates}

In this section we derive an a priori estimate for solutions of ODEs that is related to the energy estimates for PDEs that we examine in later chapters. The uniqueness theorem~\ref{thm:ode-uniqueness} is an immediate consequence of this result. To derive our estimate we need a fundamental inequality called Gronwall's inequality.

\begin{Lemma}{Gronwall's inequality}
\label{lem:gronwall}
Let
\[
\begin{aligned}u&:[a,b]\to[0,\infty),\\v&:[a,b]\to\mathbb{R},\end{aligned}
\]
be continuous functions and let $C$ be a constant. Then if
\begin{equation}\label{eq:gronwall-integral-hypothesis}
v(t)\leq C+\int_a^tv(s)u(s)\,\mathrm{d}s
\end{equation}
for $t\in[a,b]$, it follows that
\begin{equation}\label{eq:gronwall-exponential-bound}
v(t)\leq C\exp\left(\int_a^tu(s)\,\mathrm{d}s\right)
\end{equation}
for $t\in[a,b]$.
\end{Lemma}

The proof of this is left as an exercise.

\begin{Lemma}{Energy estimate for ODEs}
\label{lem:ode-energy-estimate}
Let $\bm{F}:\mathbb{R}\times\mathbb{R}^n\to\mathbb{R}^n$ satisfy the hypotheses of Theorem~\ref{thm:ode-existence}, in particular let it be uniformly Lipschitz in its second variable with Lipschitz constant $\gamma$ (cf. \eqref{eq:uniform-lipschitz-bound}). Let $\bm{y}_1$ and $\bm{y}_2$ be solutions of \eqref{eq:ode-initial-value-dynamics} on the interval $[t_0,T]$; i.e.,
\[
\bm{y}_i'(t)=\bm{F}(t,\bm{y}_i(t))
\]
for $i=1,2$ and $t\in[t_0,T]$. Then
\begin{equation}\label{eq:ode-energy-estimate}
|\bm{y}_1(t)-\bm{y}_2(t)|^2\leq|\bm{y}_1(t_0)-\bm{y}_2(t_0)|^2e^{2\gamma(t-t_0)}.
\end{equation}

\end{Lemma}

\begin{Proof}
We begin by using the differential equation, the Cauchy-Schwarz inequality and the Lipschitz condition to derive the following inequality.
\[
\begin{aligned}
&|\bm{y}_1(t)-\bm{y}_2(t)|^2\\
&=|\bm{y}_1(t_0)-\bm{y}_2(t_0)|^2+\int_{t_0}^t\dfrac{\mathrm{d}}{\mathrm{d}s}|\bm{y}_1(s)-\bm{y}_2(s)|^2\,\mathrm{d}s\\
&=|\bm{y}_1(t_0)-\bm{y}_2(t_0)|^2\\
&\quad+\int_{t_0}^t2(\bm{y}_1(s)-\bm{y}_2(s))\cdot(\bm{F}(s,\bm{y}_1(s))-\bm{F}(s,\bm{y}_2(s)))\,\mathrm{d}s\\
&\leq|\bm{y}_1(t_0)-\bm{y}_2(t_0)|^2+\int_{t_0}^t2|\bm{y}_1(s)-\bm{y}_2(s)|\,|\bm{F}(s,\bm{y}_1(s))-\bm{F}(s,\bm{y}_2(s))|\,\mathrm{d}s\\
&\leq|\bm{y}_1(t_0)-\bm{y}_2(t_0)|^2+\int_{t_0}^t2\gamma|\bm{y}_1(s)-\bm{y}_2(s)|^2\,\mathrm{d}s.
\end{aligned}
\]
Now \eqref{eq:ode-energy-estimate} follows directly from Gronwall's inequality.
\end{Proof}

Note we can derive the uniqueness result for ODEs (Theorem~\ref{thm:ode-uniqueness}) by simply setting $\bm{y}_1(t_0)=\bm{y}_2(t_0)$ and using \eqref{eq:ode-energy-estimate}. Also observe that these results are indeed obtained \emph{a priori}: nothing we did depended on the existence of a solution, only on the equations that a solution would satisfy if it \emph{did} exist.

\paragraph*{Maximum principle for analytic functions}

As an example of an \emph{a posteriori} result we consider the following theorem.

\begin{Theorem}{Maximum modulus principle}
\label{thm:maximum-modulus}
Let $D\subset\mathbb{C}$ be a bounded domain and let $f$ be analytic on $D$ and continuous on the closure of $D$. Then $|f|$ achieves its maximum on the boundary of $D$; i.e., there exists $z_0\in\partial D$ such that
\begin{equation}\label{eq:analytic-maximum-modulus}
|f(z_0)|=\sup_{z\in\overline{D}}|f(z)|.
\end{equation}

\end{Theorem}

The reader is encouraged to prove this using Cauchy's integral formula (cf. Problem~\ref{ex:maximum-modulus-from-cauchy}). Such a proof, based on an explicit representation for the function $f$, is \emph{a posteriori}. We note, however, that it is possible to give an \emph{a priori} proof of the result; and Chapter~\ref{chap:maximum-principles} is dedicated to finding \emph{a priori} maximum principles for PDEs.

\subsubsection{Smoothness}\label{subsec:smoothness}

One of the most important modern techniques for proving the existence of a solution to a partial differential equation is the following process.
\begin{enumerate}
\item Convert the original PDE into a "weak" form that might conceivably have very rough solutions.
\item Show that the weak problem has a solution.
\item Show that the solution of the weak equation actually has more smoothness than one would have at first expected.
\item Show that a "smooth" solution of the weak problem is a solution of the original problem.
\end{enumerate}

We give a preview of parts one, two, and four of this process in Section~\ref{subsec:laplace-equation} below, but in this section let us consider precursors of the methods for part three: showing smoothness.

\paragraph*{Smoothness of solutions of ODEs}

The following is an example of a "bootstrap" proof of regularity in which we use the fact that $\bm{y}\in C^0$ to show that $\bm{y}\in C^1$, etc. Note that this result can be used to prove the regularity portion of Theorem~\ref{thm:ode-existence} (which asserted the existence of a $C^1$ solution).

\begin{Theorem}
\label{thm:ode-bootstrap-regularity}
If $\bm{F}:\mathbb{R}\times\mathbb{R}^n\to\mathbb{R}^n$ is in $C^{m-1}(\mathbb{R}\times\mathbb{R}^n)$ for some integer $m\geq1$, and $\bm{y}\in C^0(\mathbb{R})$ satisfies the integral equation
\begin{equation}\label{eq:ode-bootstrap-integral-equation}
\bm{y}(t)=\bm{y}(t_0)+\int_{t_0}^t\bm{F}(s,\bm{y}(s))\,\mathrm{d}s,
\end{equation}
then in fact $\bm{y}\in C^m(\mathbb{R})$.
\end{Theorem}

\begin{Proof}
Since $\bm{F}(s,\bm{y}(s))$ is continuous, we can use the Fundamental Theorem of Calculus to deduce that the right-hand side of \eqref{eq:sturm-liouville-kernel-eigenproblem} is continuously differentiable, so the left-hand side must be as well, and
\begin{equation}\label{eq:ode-bootstrap-differential-equation}
\bm{y}'(t)=\bm{F}(t,\bm{y}(t)).
\end{equation}
Thus, $\bm{y}\in C^1(\mathbb{R})$. If $\bm{F}$ is in $C^1$, we can repeat this process by noting that the right-hand side of \eqref{eq:ode-bootstrap-differential-equation} is differentiable (so the left-hand side is as well) and
\[
\bm{y}''(t)=\bm{F}'_{\bm{y}}(t,\bm{y}(t))\cdot\bm{y}'(t)+\bm{F}_t(t,\bm{y}(t)),
\]
so $\bm{y}\in C^2(\mathbb{R})$. This can be repeated as long as we can take further continuous derivatives of $\bm{F}$. We conclude that, in general, $\bm{y}$ has one order of differentiability more than $\bm{F}$.
\end{Proof}

\paragraph*{Smoothness of analytic functions}

A stronger result can be obtained for analytic functions by using Cauchy's integral formula.

\begin{Theorem}
\label{thm:analytic-infinite-differentiability}
If a function $f:\mathbb{C}\to\mathbb{C}$ is analytic at $z_0\in\mathbb{C}$ (i.e., if it has at least one complex derivative in a neighborhood of $z_0$), then it has complex derivatives of arbitrary order. In fact,
\begin{equation}\label{eq:cauchy-higher-derivatives}
f^{(n)}(z_0)=\dfrac{n!}{2\pi i}\oint_C\dfrac{f(z)}{(z-z_0)^{n+1}}\,\mathrm{d}z
\end{equation}
for any simple, closed, positively oriented curve $C$ lying in a simply connected domain in which $f$ is analytic and having $z_0$ in its interior.
\end{Theorem}

The proof can be obtained by differentiating Cauchy's integral formula \eqref{eq:cauchy-integral-formula} under the integral sign. This is a common technique in PDEs, and one with which the reader should be familiar (cf. Problem~\ref{ex:differentiation-under-integral}).

\paragraph*{Problems}

\begin{enumerate}[label=\thesection.\arabic*.,ref=\thesection.\arabic*,series=book-problems]

\item\label{ex:picard-uniform-convergence} Let $\bm{y}_k$ be the sequence defined by \eqref{eq:picard-iteration}. Show that
\[
|\bm{y}_{k+1}(t)-\bm{y}_k(t)|\leq\gamma(t-t_0)\max_{\tau\in[t_0,t]}|\bm{y}_k(\tau)-\bm{y}_{k-1}(\tau)|.
\]
Use this to show that the sequence converges uniformly for $t_0\leq t\leq T$ for any $T<t_0+1/\gamma$.

\item\label{ex:sphere-implicit-function} Use the implicit function theorem to determine when the equation
\[
x^2+y^2+z^2=1
\]
defines implicitly a function $\hat{x}(y,z)$. Give a geometric interpretation of this result.

\item\label{ex:implicit-function-regularity} Show that if $\bm{F}$ as described in Theorem~\ref{thm:implicit-function} is $C^k$, then $\hat{\bm{y}}$ is $C^k$ as well. Hint: First consider difference quotients to show $\hat{\bm{y}}$ is $C^1$.

\item\label{ex:nonunique-minimizers} Show that there is an infinite family of minimizers of
\[
\mathcal{E}(u):=\int_0^1(1-(u'(t))^2)^2\,\mathrm{d}t
\]
over the set of all piecewise $C^1$ functions satisfying $u(0)=u(1)=0$.

\item\label{ex:nonattained-minimum} Show that there is no piecewise $C^1$ minimizer of
\[
\mathcal{E}(u):=\int_0^1u(t)^2+(1-(u'(t))^2)^2\,\mathrm{d}t
\]
satisfying $u(0)=u(1)=0$. Hint: Use a sequence of the solutions of the previous problem to show that a minimizer $\tilde{u}$ would have to satisfy $\mathcal{E}(\tilde{u})=0$. Remark: Minimization problems with features like these arise in the modeling of phase transitions.

\item\label{ex:nonuniform-initial-data-continuity} Give an example that shows that $\delta$ in Theorem~\ref{thm:initial-data-continuity} cannot be chosen independent of $t$ as $t\to\infty$.

\item\label{ex:distinct-eigenvalue-stability} Prove Theorem~\ref{thm:constant-coefficient-stability} in the case where the eigenvalues of $\bm{A}$ are distinct.

\item\label{ex:fundamental-matrix-existence} Prove the existence and uniqueness of the solution of \eqref{eq:fundamental-matrix-dynamics}, \eqref{eq:fundamental-matrix-normalization}.

\item\label{ex:gronwall-proof} Prove Gronwall's inequality.

\item\label{ex:maximum-modulus-from-cauchy} Prove Theorem~\ref{thm:maximum-modulus} using Cauchy's integral formula.

\item\label{ex:differentiation-under-integral} Suppose $g:\mathbb{R}^2\to\mathbb{R}$ is $C^1$. Define
\[
f(x):=\int_a^bg(x,y)\,\mathrm{d}y.
\]
Use difference quotients to show that one can differentiate $f$ "under the integral sign."

\end{enumerate}

\subsection{Elementary Partial Differential Equations}\label{sec:elementary-pdes}

In the last section we discussed the basic types of mathematical questions that are considered throughout the rest of this book, and we looked at how those questions had been answered by two subdisciplines of PDEs: ODEs and complex variables. We now look at how these questions are often approached in elementary courses on partial differential equations. To do this, we consider three basic PDEs (Laplace's equation, the heat equation and the wave equation). Although we sometimes use an analytical approach to investigate their character, our basic technique is the explicit calculation of solutions. At this point we are not terribly concerned with either rigor or generality but rather with foreshadowing material to come; all of the methods and observations presented here are generalized later on.
